\section{Estado del arte}
\label{sec:resultados_revision}
\subsection{Síntesis metodológica de la revisión}
El estado del arte de este TFM se apoya en una revisión sistemática de la literatura desarrollada como estudio independiente y enviada a evaluación en la revista \textit{Applied Sciences}. La revisión siguió la guía de Kitchenham \cite{kitchenham2007} y el estándar PRISMA 2020 \cite{prisma2020}, con una estrategia de búsqueda delimitada mediante el modelo PICo. La consulta se ejecutó en Scopus, Web of Science e IEEE Xplore, con restricción a publicaciones en inglés entre 2021 y 2026.

La búsqueda inicial recuperó 240 registros. Tras depuración de duplicados, cribado por título y resumen, recuperación de texto completo y evaluación de elegibilidad, el corpus final quedó constituido por 19 estudios primarios. En esta memoria se presenta la parte del estudio que resulta necesaria para tres propósitos: caracterizar el campo reciente, justificar la brecha de investigación y derivar criterios de diseño para el módulo de monitoreo planteado en este TFM.

\subsection{Hallazgos del estado del arte relevantes para el TFM}
\subsubsection{Variables de entrada y fuentes de datos}
La literatura reciente converge en un conjunto acotado de variables para evaluación de la calidad del aire en sistemas difusos. PM$_{2.5}$, CO, CO$_2$ y PM$_{10}$ aparecen con alta recurrencia en tareas de clasificación, estimación y predicción, mientras temperatura y humedad se emplean como variables de apoyo en arquitecturas interiores y urbanas \cite{ART06,ART07,A028,A135,A165,A167}. Esta pauta se observa en sistemas de monitoreo IoT con clasificación AQI/ISPU basados en PM$_{2.5}$, PM$_{10}$ y CO \cite{ART06}, en propuestas urbanas que combinan SO$_2$, NO$_2$, CO, O$_3$ y PM$_{10}$ \cite{ART07}, y en esquemas de pronóstico o ajuste donde PM$_{2.5}$ se combina con covariables microclimáticas y contaminantes adicionales \cite{A084,A135}. Para este TFM, esta evidencia delimita un núcleo de variables pertinente para diseñar un módulo urbano basado en datos abiertos y orientado a evaluación del riesgo.

La revisión también muestra una diferencia metodológica importante entre adquisición directa y reutilización de datos abiertos. Predominan los estudios sustentados en sensores propios o redes instrumentadas de bajo costo \cite{ART05,ART06,A114,A165,A167}. Los trabajos basados en datos públicos o en validación multi-localización son menos frecuentes, aunque aportan mejores condiciones de trazabilidad para contrastar generalización y comparabilidad \cite{A028,A084,A135}. Este punto es central para la propuesta del TFM, porque la procedencia del dato condiciona la reproducibilidad del módulo y la validez de su evaluación.

\subsubsection{Inferencia difusa y modelado del riesgo}
Los estudios revisados sitúan la inferencia difusa como núcleo de decisión del sistema. Las familias más recurrentes son \texttt{Mamdani} y \texttt{ANFIS}. \texttt{Mamdani} aparece asociada con mayor frecuencia a clasificación de calidad del aire, interpretación de categorías AQI y activación de alertas \cite{ART06,A165,A167}. \texttt{ANFIS} aparece en tareas de predicción, ajuste de medida y reducción de error sobre series de contaminantes o sensores de bajo costo \cite{A028,A084,A135}. Esta distinción resulta decisiva para el TFM, porque separa una lógica de diseño orientada a trazabilidad de reglas y soporte a la decisión de otra centrada en ajuste guiado por datos.

La literatura mantiene heterogeneidad en aspectos críticos del diseño difuso. Las bases de reglas, las funciones de pertenencia y el criterio de modelado del riesgo se reportan con distinto nivel de detalle. En varios estudios, la relación entre regla, clase AQI y respuesta operativa queda explícita \cite{ART05,ART06,A114,A165,A167}. En otros, el énfasis se desplaza hacia desempeño predictivo o simplificación del modelo, como ocurre en trabajos basados en \texttt{ANFIS} o reducción explicable de reglas \cite{A028,A084}. Esta diferencia afecta la auditabilidad de los resultados y limita la comparación entre sistemas que usan la misma etiqueta difusa con decisiones internas distintas. Para un módulo destinado a apoyar evaluación del riesgo y generación de alertas, esta observación refuerza la necesidad de un núcleo difuso interpretable y documentado.

\subsubsection{Plataformas de monitoreo y operación en tiempo real}
El corpus revisado confirma la presencia de una base tecnológica relativamente estable para monitoreo difuso en tiempo real o casi tiempo real. Son frecuentes las arquitecturas IoT con nodo de adquisición, procesamiento local o híbrido y publicación en nube o tablero web \cite{ART05,ART06,A114,A135,A167}. En \cite{A114}, el flujo integra sensores, inferencia, visualización y actuación sobre el ambiente mediante ESP32 y Firebase. En \cite{A135}, la red IoT alimenta un portal web para seguimiento de IAQ y pronóstico de PM$_{2.5}$. En \cite{A167}, el sistema combina sensores de bajo costo, calibración difusa y estimación de AQI en un entorno doméstico.

La operación continua no implica el mismo nivel de integración en todos los casos. Algunos trabajos reportan monitoreo y publicación de resultados con respuesta por umbral o recomendación operativa \cite{ART05,A114,A165,A167}. Otros concentran el aporte en el modelo predictivo o en la estructura de red y dejan en segundo plano la descripción del ciclo completo de adquisición, procesamiento, publicación y mantenimiento \cite{A028,ART08}. Para este TFM, la diferencia es sustantiva: el objetivo consiste en formular un modelo de inferencia y además desarrollar un módulo de monitoreo con flujo funcional completo y salida útil para alertamiento.

\subsection{Vacíos del estado del arte}
La síntesis del corpus revisado permite identificar un conjunto de vacíos que trascienden a estudios aislados y condicionan tanto la comparabilidad de resultados como el diseño de nuevas propuestas. Estos vacíos, que se presentan a continuación, afectan la procedencia del dato, la amplitud del contexto evaluado, el nivel de integración operativa y la forma en que se reporta la validación.

\begin{itemize}
\item \textbf{Generalización externa insuficiente}. La literatura presenta resultados favorables dentro del entorno de prueba de cada estudio, pero la validación en más de una ciudad, más de una red de sensores o más de una ventana temporal sigue siendo escasa \cite{ART04,A028,A084,A135}. Esta limitación ya había sido anticipada en revisiones de AQI, donde la comparabilidad entre modelos y contextos aparece como problema estructural \cite{ART01,ART02}.
\item \textbf{Comparabilidad metodológica limitada}. Los trabajos difieren en variables de entrada, tipo de salida, métrica reportada, partición de datos y nivel de detalle sobre la configuración difusa. En consecuencia, dos estudios pueden informar alto desempeño sin estar evaluando la misma tarea ni bajo un protocolo equivalente \cite{ART01,ART02,A028,A084,A135}. Esta heterogeneidad dificulta atribuir el rendimiento al diseño del sistema con la precisión necesaria para reutilizarlo en otro contexto.
\item \textbf{Cobertura restringida de variables y contextos}. El núcleo de contaminantes y covariables se concentra en configuraciones compactas, útiles para despliegue práctico, pero con menor amplitud química y contextual \cite{ART06,ART07,A135,A165,A167}. En varios casos, la evaluación se desarrolla en contextos interiores o locales, con poca discusión sobre transferencia a escenarios urbanos abiertos o a fuentes públicas heterogéneas \cite{ART05,A114,A167}.
\item \textbf{Integración operativa incompleta}. La literatura incluye ejemplos con monitoreo, tablero y respuesta automática \cite{ART05,A114,A165,A167}, pero también trabajos en los que el modelo queda planteado como componente analítico potencial sin cierre completo de la cadena operativa \cite{A028,A135}. Para una propuesta orientada a soporte de decisión, esta diferencia es crítica porque la utilidad práctica depende de la trazabilidad del flujo completo.
\item \textbf{Reporte desigual de calibración, interpretabilidad y eficiencia operativa}. Estudios recientes muestran avances claros en ajuste explicable de sensores de bajo costo \cite{A084} y en eficiencia de red para WSN urbanas \cite{ART08}. Sin embargo, esas dimensiones no se reportan de forma homogénea en el conjunto del campo. Persisten diferencias importantes en la documentación de reducción de reglas, deriva de sensores, tiempos de inferencia, consumo de recursos y desempeño de comunicación \cite{A084,ART08}. Esta dispersión limita la reproducibilidad técnica.
\end{itemize}

Estos vacíos técnicos describen un campo activo, pero todavía fragmentado en variables, diseño difuso, arquitectura y validación. Esta evidencia justifica el objetivo del TFM: desarrollar un módulo de monitoreo apoyado en datos abiertos, con inferencia difusa auditable y con criterios de evaluación comparables.

\subsection{Antecedentes comparables para esta investigación}
La Tabla~\ref{tab:antecedentes_comparables_tfm} reúne estudios del corpus que ofrecen puntos de comparación directos para esta investigación. La selección prioriza trabajos con afinidad en cuatro dimensiones: uso de datos trazables, núcleo difuso explícito, monitoreo en tiempo real o casi tiempo real, y generación de alertas o soporte a decisión.

Dentro de ese conjunto aparecen antecedentes centrados en predicción AQI urbana con datos públicos \cite{A028}, ajuste explicable de PM$_{2.5}$ sobre múltiples localizaciones \cite{A084}, control ambiental industrial con integración IoT \cite{A114}, pronóstico interior con portal web \cite{A135}, alertamiento local basado en reglas \cite{A165} y monitoreo automático indoor con FIS tipo 1 \cite{A167}. Esta combinación permite contrastar la propuesta con estudios que cubren distintos puntos de la cadena, desde el modelado hasta la operación de plataforma, sin perder afinidad con el problema de monitoreo y decisión.

\begin{longtblr}[
  label={tab:antecedentes_comparables_tfm},
  caption={Antecedentes comparables derivados del SLR para contrastar la propuesta del TFM.}
]{
  colspec={X[0.9,l,m] X[1.4,l,m] X[1.4,l,m] X[1.5,l,m] X[1.1,l,m] X[1.7,l,m]},
  width=\linewidth
}
\toprule
Estudio & Objetivo y contexto & Entradas y fuente de datos & Enfoque difuso y plataforma & Validación reportada & Límite relevante para contrastar con este TFM \\
\midrule
\cite{A028} & Predicción de AQI en contexto urbano. & Datos ambientales y de contaminantes en tiempo real; el artículo destaca NO$_2$ y CO como variables centrales. & Arquitectura híbrida con \texttt{ANFIS} dentro de un ensamble de aprendizaje automático. & Reporta \(R^2=96\%\) en particiones aleatorizadas. & El trabajo plantea integración potencial en estaciones de monitoreo, pero no documenta un despliegue operativo completo. \\
\cite{A084} & Ajuste de mediciones de PM$_{2.5}$ de sensores de bajo costo en múltiples localizaciones. & PM$_{2.5}$ y datos de co-localización con validación geográfica múltiple; publica software y datos en Zenodo. & \texttt{ANFIS} con reducción explicable de reglas; implementación en MATLAB. & Mantiene correlaciones de prueba entre 0.73 y 0.96 tras la poda de reglas. & Su foco está en calibración y ajuste de medida; no aborda alertamiento ni operación urbana de extremo a extremo. \\
\cite{A114} & Control de calidad ambiental en fábrica inteligente. & Sensores ambientales conectados a ESP32-S y envío a Firebase. & Reglas difusas más aprendizaje por refuerzo profundo; monitoreo remoto por app/web. & Precisión de 91.16\%, pérdida de 1.64\% y tiempo de inferencia de 3 ms. & El dominio es industrial interior; la solución prioriza control local y eficiencia energética. \\
\cite{A135} & Pronóstico de PM$_{2.5}$ interior con red IoT y portal web. & PM$_{2.5}$, CO$_2$, CO, temperatura y humedad; combina datos reales de distintas localizaciones y conjuntos de referencia. & \texttt{ADFIST} sobre red IoT y publicación en portal \textit{Vayuveda}. & Contraste con múltiples conjuntos de datos y validación adicional frente a \textit{benchmarks}. & El objetivo es pronóstico interior; la comparación interurbana abierta sigue siendo limitada. \\
\cite{A165} & Monitoreo de calidad del aire con nodo inteligente y alerta temprana. & CO, CO$_2$ y CH$_4$ en pruebas interiores y exteriores; transmisión por LoRa. & Fuzzy Tsukamoto en nodo sensor con Raspberry Pi como agregador. & Exactitud de 92.4\% para la clasificación del aire. & Trabaja con un conjunto reducido de variables y una semántica de alerta local. \\
\cite{A167} & Monitoreo automático de calidad del aire interior en hogar inteligente. & CO, PM, temperatura y humedad con sensores de bajo costo en entorno doméstico. & Sistema IoT de dos capas con FIS tipo 1 para calibración y estimación de AQI. & Ensayos experimentales en vivienda; el artículo reporta efectividad del sistema implementado. & El estudio se centra en entorno interior y no desarrolla validación externa entre ciudades o fuentes abiertas. \\
\bottomrule
\end{longtblr}


\subsection{Implicaciones para el diseño de la propuesta}
Los hallazgos del estado del arte delimitan decisiones de diseño que resultan difíciles de omitir en una propuesta orientada a monitoreo urbano, trazabilidad operativa y soporte a la decisión. A partir de las limitaciones detectadas en validación, comparabilidad, explicitud del núcleo difuso e integración de plataforma \cite{ART01,ART02,A084,A114,A135,A167}, la propuesta adopta los siguientes criterios:
\begin{enumerate}
\item La fuente de datos debe ser trazable y reutilizable, porque la literatura muestra una dependencia excesiva de instrumentación local con validación acotada \cite{A028,A084,A135}.
\item El componente difuso debe mantener reglas, funciones de pertenencia y criterios de salida explícitos, dado que la auditabilidad sigue siendo un punto débil en parte del campo \cite{ART05,ART06,A084,A167}.
\item La propuesta debe tratar el monitoreo como un flujo operativo completo, con ingesta, procesamiento, inferencia, publicación y alertamiento documentados \cite{A114,A165,A167}.
\item La validación debe reportarse con métricas compatibles con la salida objetivo y con una partición de datos claramente declarada \cite{ART01,ART02,A028,A135}.
\item La solución debe incorporar control de calidad de datos y una estrategia de evaluación que permita observar estabilidad fuera de un único contexto de prueba \cite{A084,ART08}.
\end{enumerate}


